\documentclass[a4paper,12pt]{article}

\input{../../Preambulo.tex}
\newcommand{\assunto}{CONJUNTOS E FUNÇÕES}
\newcommand{\atividade}{atividade} % prova ou lista ou as ou ad ou atividade
\newcommand{\NumQObj}{6}
\newcommand{\ValorProva}{2}
\pgfmathparse{\ValorProva/\NumQObj}
\edef\ValorQObj{\pgfmathresult}
%\newcommand{\ValorQObj}{0.3}
%\pgfmathparse{int(round(\NumQObj * \ValorQObj))}
%\edef\ValorProva{\pgfmathresult}
\newcommand{\tipo}{dupla} % dupla ou individual
\newcommand{\calculadora}{nao} % sim ou nao
\newcommand{\turma}{ELETROELETRÔNICA 1.\textordmasculine\ ANO}
\newcommand{\numProva}{A}
\newcommand{\professor}{Igor Martins Silva}
\newcommand{\bimestre}{3}
\newcommand{\ano}{2026}
\newcommand{\mesTexto}{agosto}
\newcommand{\mesNum}{08}
\newcommand{\dia}{27}
\newcommand{\dataTexto}{\dia\ de \mesTexto\ de \ano}
\newcommand{\dataNun}{\dia/\mesNum/\ano}

\newcommand{\titulo}{}

\ifdefstring{\atividade}{prova}
  {\renewcommand{\titulo}{Avaliação de Matemática}}
{
\ifdefstring{\atividade}{lista}
  {\renewcommand{\titulo}{Lista de exercícios de Matemática}}
{
\ifdefstring{\atividade}{as}
  {\renewcommand{\titulo}{Avaliação Somativa}}
{
\ifdefstring{\atividade}{ad}
  {\renewcommand{\titulo}{Avaliação Diagnóstica}}
  {\renewcommand{\titulo}{Atividade de Matemática}}
}}}

\newcommand{\emtipo}{%
  \ifdefstring{\tipo}{individual}{\tipo}{em \tipo}%
}

\edef\pathCapa{\currfiledir}

\newcommand{\subtitulobase}[3]{%
    % #1 = título da 1ª coluna
    % #2 = conteúdo da 1ª coluna
    % #3 = mostra número da prova? (vazio ou algo)
    
    \noindent
    \begin{minipage}{2.8cm}
        \includegraphics[width=2.8cm]{\pathCapa Logo_Cefet.png}
    \end{minipage}
    \hfill
    \begin{minipage}{\dimexpr\textwidth-6.2cm\relax}
        \begin{center}
        CEFET-MG -- Campus Contagem \\[3pt]
        \titulo\ -- \bimestre.\textordmasculine\ bimestre de \ano \\[3pt]
        \professor
    \end{center}
    \end{minipage}
    \hfill
    \begin{minipage}{3.2cm}
        \includegraphics[width=3.2cm]{\pathCapa Brasao_Brasil.png}
    \end{minipage}
    
    \begin{center}
        \vspace{15pt}
        \begin{tabular}{|C{210pt}|C{70pt}|C{210pt}|}
            \hline
            \textbf{#1} &
            \textbf{DATA} &
            \textbf{TURMA} \\
            \hline
            #2 & \dataNun & \turma \\
            \hline
        \end{tabular}
    \end{center}
    
    #3
    
    \vspace{15pt}
}

\newcommand{\subtitulolis}{%
  \subtitulobase{ASSUNTO}{\assunto}{}%
}

\newcommand{\subtitulopr}{%
  \subtitulobase{ASSUNTO}{\assunto}{\boxed{\numProva}}%
}

\newcommand{\subtituloas}{%
  \subtitulobase{DISCIPLINA}{MATEMÁTICA}{\boxed{\numProva}}%
}

\newcommand{\subtituload}{%
  \subtitulobase{DISCIPLINA}{MATEMÁTICA}{}%
}

\newcommand{\valorquestao}{}

\newcommand{\definevalorquestao}{%
    \ifdefstring{\atividade}{lista}{%
        % Não faz nada para lista
    }{%
        \ifdefstring{\modo}{objetivo}{%
            \renewcommand{\valorquestao}{\ValorQObj ponto}
        }{%
            \renewcommand{\valorquestao}{\ValorQDisc pontos}
        }%
    }%
}

\newenvironment{exercicioBanco}[1][]{%
    \ifdefstring{\atividade}{lista}{%
        \begin{exercicio}%
    }{%
        \begin{exercicio}[#1]%
    }%
}{%
    \end{exercicio}
}

\ifdefstring{\atividade}{lista}{
    \pagecolor{background-color}
}{
    
}



\hypersetup{
    pdfauthor={\professor},
    pdftitle={\titulo \-- \assunto \-- \turma \-- \numProva},
}

\title{\titulo \-- \assunto \-- \turma \-- \numProva}
\author{\professor}
\date{\data}

%************* INÍCIO DO DOCUMENTO *************
\begin{document}
    \NoBgThispage
    \pagenumbering{alph}
    %
        
    \phantomsection
    \addcontentsline{toc}{section}{Capa}
    \input{../../AT_Capa_Gabarito}
    \newpage
    \null
    \thispagestyle{empty}
    \addtocounter{page}{1}
    \newpage
        
    \NoBgThispage
    \pagenumbering{arabic}
    %
        
    \subtitulopr
    
    \vspace{15pt}
    \newcommand{\modo}{objetivo} % objetivo ou discursivo
    
    \phantomsection
    \addcontentsline{toc}{section}{Exercício 1}
    %%%%%%%%%%%%%%%%%%%%%%%%%%%%%%%%%
%
%       EXERCÍCIO
%
%%%%%%%%%%%%%%%%%%%%%%%%%%%%%%%%%

\ifdefstring{\atividade}{prova}{%
    \ifdefstring{\modo}{objetivo}{%
        \renewcommand{\valorquestao}{\ValorQObj\ ponto}
    }{%
        \renewcommand{\valorquestao}{\ValorQDisc\ pontos}
    }%
}%

\begin{exercicioBanco}[\valorquestao]
Em uma turma, \(40\) estudantes responderam a uma pesquisa sobre duas plataformas de streaming, \(A\) e \(B\).
Sabe-se que:

\begin{enumerate}[\(\bullet\)]
    \begin{multicols}{2}
        \item \(18\) estudantes utilizam a plataforma \(A\);
        \item \(9\) estudantes utilizam ambas as plataformas;
        \item \(11\) estudantes não utilizam nenhuma das duas plataformas.
    \end{multicols}
\end{enumerate}

Sabendo que todos os \(40\) estudantes responderam à pesquisa, o número de estudantes que utilizam a plataforma \(B\) é:

% Define as alternativas
\newcommand{\alternativas}{%
    \begin{center}
        \begin{tabularx}{\textwidth}{XXXXX}
            (a) \(11\). &
            (b) \(14\). &
            (c) \(17\). &
            (d) \(20\). &
            (e) \(22\).
        \end{tabularx}
    \end{center}
}

% Define a resposta correta
\newcommand{\resposta}{D}

% Lógica condicional para exibição
\ifdefstring{\atividade}{lista}{%
    \alternativas
    \vspace{0.5em}
    
    \noindent\textbf{Resposta:} letra \textbf{\resposta}.
}{%
    \ifdefstring{\modo}{objetivo}{%
        \alternativas
    }{%
        % modo = discursiva → não mostra alternativas
    }
}
\end{exercicioBanco}


    \vspace{15pt}
    
    \phantomsection
    \addcontentsline{toc}{section}{Exercício 2}
    \input{../C_Exercicio5_2}
    \vspace{15pt}
    
    \phantomsection
    \addcontentsline{toc}{section}{Exercício 3}
    %%%%%%%%%%%%%%%%%%%%%%%%%%%%%%%%%
%
%       EXERCÍCIO
%
%%%%%%%%%%%%%%%%%%%%%%%%%%%%%%%%%

\ifdefstring{\atividade}{prova}{%
    \ifdefstring{\modo}{objetivo}{%
        \renewcommand{\valorquestao}{\ValorQObj\ ponto}
    }{%
        \renewcommand{\valorquestao}{\ValorQDisc\ pontos}
    }%
}%

\begin{exercicioBanco}[\valorquestao]
Dados os conjuntos \(A = [0,7[\), \(B = ]-3,3]\) e \(E = ]-1,5[\), calcule \(E\setminus(A\cap B)\).

% Define as alternativas
\newcommand{\alternativas}{%
    \begin{center}
        \begin{tabularx}{\textwidth}{XXX}
            (a) \(]-1,0[ \cup ]3,5[\). &
            (b) \(]-1,0] \cup [3,5[\). &
            (c) \(]-1,3]\). \\[5pt]
            (d) \([0,3]\). &
            (e) \(]-1,5[ \setminus [0,3[\).
        \end{tabularx}
    \end{center}
}

% Define a resposta correta
\newcommand{\resposta}{A}

% Lógica condicional para exibição
\ifdefstring{\atividade}{lista}{%
    \alternativas
    \vspace{0.5em}
    
    \noindent\textbf{Resposta:} letra \textbf{\resposta}.
}{%
    \ifdefstring{\modo}{objetivo}{%
        \alternativas
    }{%
        % modo = discursiva → não mostra alternativas
    }
}
\end{exercicioBanco}


    \vspace{15pt}
    
    \phantomsection
    \addcontentsline{toc}{section}{Exercício 4}
    %%%%%%%%%%%%%%%%%%%%%%%%%%%%%%%%%
%
%       EXERCÍCIO
%
%%%%%%%%%%%%%%%%%%%%%%%%%%%%%%%%%

\ifdefstring{\atividade}{prova}{%
    \ifdefstring{\modo}{objetivo}{%
        \renewcommand{\valorquestao}{\ValorQObj\ ponto}
    }{%
        \renewcommand{\valorquestao}{\ValorQDisc\ pontos}
    }%
}%

\begin{exercicioBanco}[\valorquestao]
Ordene os números racionais \(p = \dfrac{7}{12}\), \(q = \dfrac{3}{4}\) e \(r = \dfrac{5}{6}\).

% Define as alternativas
\newcommand{\alternativas}{%
    \begin{center}
        \begin{tabularx}{\textwidth}{XXXXX}
            (a) \(p < q < r\). &
            (b) \(q < p < r\). &
            (c) \(r < q < p\). &
            (d) \(q < r < p\). &
            (e) \(r < p < q\).
        \end{tabularx}
    \end{center}
}

% Define a resposta correta
\newcommand{\resposta}{A}

% Lógica condicional para exibição
\ifdefstring{\atividade}{lista}{%
    \alternativas
    \vspace{0.5em}
    
    \noindent\textbf{Resposta:} letra \textbf{\resposta}.
}{%
    \ifdefstring{\modo}{objetivo}{%
        \alternativas
    }{%
        % modo = discursiva → não mostra alternativas
    }
}
\end{exercicioBanco}


    \vspace{15pt}
    
    \phantomsection
    \addcontentsline{toc}{section}{Exercício 5}
    %%%%%%%%%%%%%%%%%%%%%%%%%%%%%%%%%
%
%       EXERCÍCIO
%
%%%%%%%%%%%%%%%%%%%%%%%%%%%%%%%%%

\ifdefstring{\atividade}{prova}{%
    \ifdefstring{\modo}{objetivo}{%
        \renewcommand{\valorquestao}{\ValorQObj\ ponto}
    }{%
        \renewcommand{\valorquestao}{\ValorQDisc\ pontos}
    }%
}%

\begin{exercicioBanco}[\valorquestao]
Seja \(f: \bbR \to \bbR\) uma função tal que \(f(x + y) = f(x) + f(y)\), para todo \(x, y \in \bbR\), e \(f(1) = 2\). Determine o valor de \(f(4)\).

% Define as alternativas
\newcommand{\alternativas}{%
    \begin{center}
        \begin{tabularx}{\textwidth}{XXXXX}
            (a) \(4\). &
            (b) \(6\). &
            (c) \(8\). &
            (d) \(10\). &
            (e) \(16\).
        \end{tabularx}
    \end{center}
}

% Define a resposta correta
\newcommand{\resposta}{C}

% Lógica condicional para exibição
\ifdefstring{\atividade}{lista}{%
    \alternativas
    \vspace{0.5em}
    
    \noindent\textbf{Resposta:} letra \textbf{\resposta}.
}{%
    \ifdefstring{\modo}{objetivo}{%
        \alternativas
    }{%
        % modo = discursiva → não mostra alternativas
    }
}
\end{exercicioBanco}


    \vspace{15pt}
    
    \newpage 
    
    \phantomsection
    \addcontentsline{toc}{section}{Exercício 6}
    %%%%%%%%%%%%%%%%%%%%%%%%%%%%%%%%%
%
%       EXERCÍCIO
%
%%%%%%%%%%%%%%%%%%%%%%%%%%%%%%%%%

\ifdefstring{\atividade}{prova}{%
    \ifdefstring{\modo}{objetivo}{%
        \renewcommand{\valorquestao}{\ValorQObj\ ponto}
    }{%
        \renewcommand{\valorquestao}{\ValorQDisc\ pontos}
    }%
}%

\begin{exercicioBanco}[\valorquestao]
Considere a função $f:[-4,5] \to [-2,4]$, cujo gráfico está abaixo.

\begin{center}
    \begin{tikzpicture}[scale=0.45]
        % Eixos
        \draw[->] (-5,0) -- (6,0);
        \draw[->] (0,-3) -- (0,5);

        % Marcas
        \foreach \x in {-4,-2,-1,1,2,3,4,5}
            \draw (\x,0.15) -- (\x,-0.15);

        \foreach \y in {-2,-1,1,2,3,4}
            \draw (0.15,\y) -- (-0.15,\y);

        % Rótulos principais
        \node[above] at (-4,0) {$-4$};
        \node[above] at (1,0) {$1$};
        \node[above] at (5,0) {$5$};
        \node[left] at (0,4) {$4$};
        \node[right] at (0,-2) {$-2$};
        \node[left] at (0,-1) {$-1$};

        % Segmento inclinado
        \draw[thick] (-4,-2) -- (1,4);

        % Pontos fechados
        \fill (-4,-2) circle (4pt);
        \fill (1,4) circle (4pt);

        % Segmento horizontal inferior
        \draw[thick] (1,-1) -- (5,-1);

        % Círculo aberto em (1,-1)
        \draw[thick] (1,-1) circle (4pt);

        % Ponto fechado em (5,-1)
        \fill (5,-1) circle (4pt);

        % Linhas auxiliares
        \draw[dashed] (1,1.2) -- (1,4);
        \draw[dashed] (0,-1) -- (1,-1) -- (1,0);
        \draw[dashed] (5,0) -- (5,-1);
        \draw[dashed] (0,4) -- (1,4);
        \draw[dashed] (0,-2) -- (-4,-2) -- (-4,0);
    \end{tikzpicture}
\end{center}

Determine o valor de \(\dfrac{f(4)+f(-4)}{f(1)-f(5)}\).

% Define as alternativas
\newcommand{\alternativas}{%
    \begin{center}
        \begin{tabularx}{\textwidth}{XXXXX}
            (a) \(-1\). &
            (b) \(\dfrac{3}{5}\). &
            (c) \(-\dfrac{5}{3}\). &
            (d) \(1\). &
            (e) \(-\dfrac{3}{5}\).
        \end{tabularx}
    \end{center}
}

% Define a resposta correta
\newcommand{\resposta}{E}

% Lógica condicional para exibição
\ifdefstring{\atividade}{lista}{%
    \alternativas
    \vspace{0.5em}
    
    \noindent\textbf{Resposta:} letra \textbf{\resposta}.
}{%
    \ifdefstring{\modo}{objetivo}{%
        \alternativas
    }{%
        % modo = discursiva → não mostra alternativas
    }
}
\end{exercicioBanco}


    \vspace{15pt}
    
    \newpage
    \phantomsection
    \addcontentsline{toc}{section}{Rascunho}
    \input{../../Rascunho}
    
\end{document}
%*************** FINAL DO DOCUMENTO ***************
